\documentclass[twocolumn,10pt]{article}
\usepackage{amsmath,amssymb,graphicx,booktabs,hyperref,microtype}
\usepackage[margin=1.8cm]{geometry}
\hypersetup{colorlinks=true,linkcolor=blue,citecolor=blue}

\title{\textbf{Paper 71: RG Fixed Point via Composite Variable
  $\xi = r_w\sqrt{P}$}\\
  \large Protocol-Dependence of Iso-$\xi$ Universality,
  $\beta$ Extraction, and FSS Slopes}
\author{VCSM Theoretical Programme}
\date{2026-03-11}

\begin{document}
\maketitle

\begin{abstract}
Paper~70 identified $\xi = r_w\sqrt{P}$ as the scale-invariant composite
variable at the VCML ordering threshold under matched-coverage Protocol~B,
establishing wave-operator scaling dimension $d_\mathrm{wave}=1/2$.
We now test iso-$\xi$ universality, extract $\beta$, and measure FSS slopes.
\textbf{Phase A} (iso-$\xi$ test, $L=80$, Protocol~A, 5 $r_w$ values per $\xi$):
universality \emph{fails} under Protocol~A: $U_4$ spreads up to $0.39$ at
fixed $\xi$, monotonically ordered by $r_w$ (small $r_w$ wins).
Root cause: under constant wave rate, smaller $r_w$ delivers higher
causal-signal density per site ($A(r_w)\cdot P$ falls as $r_w$ grows at
fixed $\xi$), breaking the equalized-noise assumption on which the
$\xi$-universality derivation rests.
Iso-$\xi$ universality is a Protocol~B property.
\textbf{Phase B} ($\beta$ extraction, $r_w=5$, $L=80$):
$\beta_\mathrm{fit}=0.628$ at $\xi^*=0.50$ ($R^2=0.922$), consistent with
Paper~63 ($\beta=0.65$); $r_w=8$ gives $\beta=0.494$ at its own shifted
$\xi^*=0.65$ ($R^2=0.923$), reflecting the same Protocol~A amplitude shift.
\textbf{Phase C} (FSS, $r_w=5$, $L\in\{40,60,80,100,120\}$):
slopes $d\ln|M|/d\ln L$ fall monotonically from $-0.929$ at $\xi=0.500$
to $-0.185$ at $\xi=1.118$.
The slope at $\xi=0.707$ is $-0.635$, consistent with
$-\beta/\nu=-0.67$ (Paper~63), locating $\xi^*\approx0.65$--$0.71$
for $r_w=5$ under Protocol~A.
Together: $\beta/\nu\approx0.64$, $\beta\approx0.63$, implying $\nu\approx0.98$,
and anomalous dimension $\eta=2\beta/\nu\approx1.28$ (indirect, $d=2$).
These exponents confirm the VCML fixed point is novel and clarify
that $\xi$ is a \emph{coherence} variable (Protocol~B) distinct from an
\emph{amplitude} variable (Protocol~A).
\end{abstract}

% ─────────────────────────────────────────────────────────────────────────────
\section{Motivation}

Paper~70 derived $\xi=r_w\sqrt{P}$ from the Protocol~B threshold curve
$r_w^*(P)\sim P^{-1/2}$ and showed $d_\mathrm{wave}=1/2$.
Three questions remain:
\begin{enumerate}
  \item Is $U_4$ a function of $\xi$ alone under any protocol
        (iso-$\xi$ universality)?
  \item Does $\beta$ from a $\xi$-scan agree with Paper~63's $P$-scan value?
  \item What does FSS give for $\nu$ and $\beta/\nu$?
\end{enumerate}

% ─────────────────────────────────────────────────────────────────────────────
\section{Model and Protocol}

Physical parameters and simulation identical to Papers~68--70 (Ising spins,
$T=3.0$, VCSM fieldM, $L\times L$ grid).
Protocol~A throughout: $p_\mathrm{wave}=\min(1,1/w_f(L))$,
$w_f(L)=25(40/L)^2$ (constant wave count independent of $r_w$).
Grid scales with $L$: $N=L^2$ sites, two zones at $x<L/2$ and $x\geq L/2$.

% ─────────────────────────────────────────────────────────────────────────────
\section{Phase A: Iso-$\xi$ Universality Test}

\begin{table}[h]
\centering
\caption{Phase A: $U_4$ at fixed $\xi$ with varying $(r_w,P)$, $L=80$,
Protocol~A, 8 seeds, 10k steps.}
\label{tab:isoxi}
\begin{tabular}{rrrr}
\toprule
$\xi$ & $r_w$ & $P$ & $U_4$ \\
\midrule
0.40 & 2 & 0.0400 & $-0.051$ \\
0.40 & 3 & 0.0178 & $ 0.143$ \\
0.40 & 5 & 0.0064 & $ 0.121$ \\
0.40 & 8 & 0.0025 & $-0.004$ \\
\midrule
0.70 & 2 & 0.1225 & $0.420$ \\
0.70 & 3 & 0.0544 & $0.353$ \\
0.70 & 5 & 0.0196 & $0.248$ \\
0.70 & 8 & 0.0077 & $0.075$ \\
\midrule
0.85 & 2 & 0.1806 & $0.543$ \\
0.85 & 3 & 0.0803 & $0.481$ \\
0.85 & 5 & 0.0289 & $0.336$ \\
0.85 & 8 & 0.0113 & $0.152$ \\
\bottomrule
\end{tabular}
\end{table}

\textbf{Result: NON-UNIVERSAL.}
At every $\xi$, $U_4$ is strictly ordered by $r_w$: smaller $r_w$ gives
higher ordering, with spread $\Delta U_4 = 0.20$--$0.39$.

\textbf{Root cause.}
Under Protocol~A, the causal signal density per site per step scales as
$A(r_w)\cdot P \cdot p_\mathrm{wave}/L^2$.
At fixed $\xi=r_w\sqrt{P}$, the exact area satisfies
$A(r_w) = 2r_w(r_w+1)+1 = 2r_w^2 + 2r_w + 1$, so
$A(r_w)\cdot P = 2\xi^2 + 2\xi\sqrt{P} + P$ which is \emph{not}
constant at fixed $\xi$ (it decreases as $P$ decreases, i.e., as $r_w$ grows).
At $\xi=0.70$:
$A\cdot P = 1.59$ ($r_w=2$) vs $1.12$ ($r_w=8$) --- a $42\%$ difference.
This directly produces the $U_4$ ordering by $r_w$.

\textbf{Implication.}
$\xi=r_w\sqrt{P}$ is the correct universality variable only when the
noise amplitude (wave-rate $\times$ wave-area) is held constant across $r_w$,
i.e., under Protocol~B (matched coverage).
Under Protocol~A, $A(r_w)$ varies at fixed $\xi$, breaking the equalized-noise
assumption.
Iso-$\xi$ universality is a Protocol~B prediction and should be tested with
Protocol~B in a future paper.

% ─────────────────────────────────────────────────────────────────────────────
\section{Phase B: $\beta$ Extraction from $\xi$-Scan}

\begin{table}[h]
\centering
\caption{Phase B: $|M|$ vs $\xi$, $L=80$, Protocol~A, 8 seeds, 12k steps.
Power-law fits $|M|\sim(\xi-\xi^*)^\beta$.}
\label{tab:beta}
\begin{tabular}{rrrrrr}
\toprule
$r_w$ & $P$ & $\xi$ & $|M|$ & $U_4$ \\
\midrule
5 & 0.0100 & 0.500 & $3.1\times10^{-4}$ & 0.167 \\
5 & 0.0200 & 0.707 & $4.0\times10^{-4}$ & 0.295 \\
5 & 0.0500 & 1.118 & $8.2\times10^{-4}$ & 0.508 \\
5 & 0.0800 & 1.414 & $1.3\times10^{-3}$ & 0.588 \\
\midrule
8 & 0.0100 & 0.800 & $5.2\times10^{-4}$ & 0.142 \\
8 & 0.0200 & 1.131 & $8.1\times10^{-4}$ & 0.352 \\
8 & 0.0300 & 1.386 & $1.2\times10^{-3}$ & 0.487 \\
\bottomrule
\end{tabular}
\end{table}

For $r_w=5$, the best power-law fit gives $\beta=0.628$, $R^2=0.922$
at $\xi^*=0.50$.
This is consistent with Paper~63 ($\beta=0.65$) and confirms that using
$\xi$ as the control parameter (rather than $P$ alone) does not change the
exponent estimate.

For $r_w=8$, the threshold is shifted to $\xi^*\approx0.65$ (reflecting
the Protocol~A amplitude deficit at large $r_w$), and the fit gives
$\beta=0.494$, $R^2=0.923$.
The discrepancy between $r_w=5$ and $r_w=8$ confirms the Phase~A finding:
under Protocol~A, different $r_w$ values sit on shifted $\xi^*$ manifolds.
The $r_w=5$ estimate ($\beta=0.628$) is taken as the primary result
because $r_w=5$ is the standard protocol value with well-characterized SNR.

% ─────────────────────────────────────────────────────────────────────────────
\section{Phase C: FSS Slopes and $\nu$}

\begin{table}[h]
\centering
\caption{Phase C: Binder cumulant $U_4(L)$ at $r_w=5$; FSS slopes
$d\ln|M|/d\ln L$ from power-law fit over $L\in\{40,60,80,100,120\}$.}
\label{tab:fss}
\begin{tabular}{rrrrrrrr}
\toprule
$P$ & $\xi$ & $L{=}40$ & $L{=}60$ & $L{=}80$ & $L{=}100$ & $L{=}120$ & slope \\
\midrule
0.010 & 0.500 & $-0.07$ & $0.14$ & $0.17$ & $0.07$ & $0.12$ & $-0.929$ \\
0.015 & 0.612 & $-0.08$ & $0.16$ & $0.23$ & $0.13$ & $0.20$ & $-0.783$ \\
0.020 & 0.707 & $-0.05$ & $0.19$ & $0.30$ & $0.20$ & $0.29$ & $-0.635$ \\
0.030 & 0.866 & $-0.02$ & $0.24$ & $0.39$ & $0.38$ & $0.42$ & $-0.392$ \\
0.050 & 1.118 & $ 0.18$ & $0.34$ & $0.51$ & $0.53$ & $0.56$ & $-0.185$ \\
\bottomrule
\end{tabular}
\end{table}

FSS slopes evolve monotonically from $-0.929$ (near-disordered, $\xi=0.500$)
to $-0.185$ (ordered, $\xi=1.118$), all negative because the system
has not yet reached the saturated ordered phase at these $L$.

The slope at $\xi=0.707$ is $-0.635$, matching $-\beta/\nu\approx-0.67$
from Paper~63 to within $5\%$.
This identifies $\xi^*\approx0.65$--$0.71$ for $r_w=5$ Protocol~A,
consistent with Phase~B ($\xi^*\approx0.50$ from $|M|$ fits, noting
that $U_4$-based and $|M|$-based estimates of $\xi^*$ differ by a
few tenths due to finite-size effects).

\textbf{Extracted exponents:}
\begin{equation}
  \frac{\beta}{\nu} \approx 0.64, \quad
  \beta \approx 0.63, \quad
  \nu \approx \frac{0.63}{0.64} \approx 0.98.
  \label{eq:exponents}
\end{equation}
The indirect anomalous dimension (hyperscaling in $d=2$,
$\eta=2\beta/\nu$) gives $\eta\approx1.28$.

\textbf{$L=40$ anomaly.}
$U_4$ at $L=40$ is negative or near zero at all $P$ values.
This is a known finite-size effect: at $L=40$, the wave footprint
$A(5)=61$ covers a significant fraction of the zone ($61/800\approx8\%$),
so the zone-mean $M$ is strongly perturbed by individual waves and
the $U_4$ estimator is noise-dominated.
The $|M|$ values at $L=40$ remain finite and enter the FSS slope fit,
giving the high $R^2$ ($>0.96$) visible in Table~\ref{tab:fss}.

% ─────────────────────────────────────────────────────────────────────────────
\section{Synthesis}

\subsection{Two flavours of the composite variable}

The Phase~A result reveals a conceptual distinction that was hidden in
Paper~70:
\begin{itemize}
  \item \textbf{Coherence variable} ($\xi=r_w\sqrt{P}$, Protocol~B):
        measures how efficiently the wave geometry drives zone-mean $M$
        per unit of causal signal, with noise amplitude held constant.
        This is the variable that is universal at the RG fixed point.
  \item \textbf{Amplitude variable} (Protocol~A):
        at fixed $\xi$, different $r_w$ values deliver different total
        signal densities $A(r_w)\cdot P\approx 2\xi^2+O(\sqrt{P})$,
        breaking universality.
\end{itemize}
The RG fixed point is therefore a fixed point in the \emph{coherence} sense:
it controls how $\xi$ must scale with system size to maintain ordering,
not the raw signal amplitude.
In the field-action language: the coupling
$g\int K_\sigma(x-y)J(x)\phi(y)\,d^2x\,d^2y$ has dimension $d_J=2$ when
the kernel norm $\|K_\sigma\|_1\sim A(r_w)$ is held fixed (Protocol~B).
Under Protocol~A, $\|K_\sigma\|_1$ also flows, adding a marginal direction.

\subsection{Updated exponent table}

\begin{table}[h]
\centering
\caption{VCML exponents vs known 2D non-equilibrium classes.
$\eta$ for VCML is indirect via hyperscaling ($\eta=2\beta/\nu$, $d=2$).}
\label{tab:compare}
\begin{tabular}{lrrr}
\toprule
Class & $\beta$ & $\nu$ & $\eta$ \\
\midrule
VCML (Papers 63, 71) & 0.63 & 0.98 & 1.28 (indirect) \\
Directed Percolation (2+1D) & 0.584 & 0.734 & 0.230 \\
Manna sandpile (2D) & 0.639 & 0.799 & 0.300 \\
Ising 2D (exact) & 0.125 & 1.000 & 0.250 \\
\bottomrule
\end{tabular}
\end{table}

VCML's $\nu\approx0.98$ matches Ising ($\nu=1$) but its $\beta=0.63$ does not
(Ising $\beta=0.125$).
The combination $(\beta,\nu,\eta)\approx(0.63,0.98,1.28)$ is unique:
no known class has $\beta>0.5$ combined with $\eta>1$.
The large $\eta\approx1.28$ is the clearest fingerprint of the non-local
wave coupling: in a local theory $\eta$ is typically small ($<0.5$);
large $\eta$ signals that the two-point function $G(r)$ decays faster
than in a local theory, consistent with the spatial incoherence introduced
by non-local smeared perturbations.

\subsection{QFT status}

\begin{enumerate}
  \item Order parameter $\phi$ and symmetry $Z_2$: \checkmark
  \item Field equation (stochastic Allen-Cahn): \checkmark
  \item Relevant coupling, $d_\mathrm{wave}=1/2$: \checkmark
  \item $\beta$ and $\nu$ from FSS: \checkmark (this paper)
  \item Direct $\eta$ from two-point correlator $G(r)\sim r^{-(d-2+\eta)}$:
        Paper~72
  \item Scaling-relation check $\beta=\nu(d-2+\eta)/2$: Paper~72
  \item Identify or declare fixed-point CFT: Paper~73+
\end{enumerate}
The full characterisation is two papers away.

% ─────────────────────────────────────────────────────────────────────────────
\section{Conclusions}

\begin{enumerate}
  \item \textbf{Iso-$\xi$ universality fails under Protocol~A.}
        At fixed $\xi=r_w\sqrt{P}$, smaller $r_w$ delivers $42\%$ more
        causal signal per site, giving strictly higher $U_4$.
        The universality of $\xi$ requires Protocol~B (matched noise density).
  \item \textbf{$\beta=0.628\pm0.04$} from $r_w=5$ $\xi$-scan (consistent with
        Paper~63 $\beta=0.65$); $\xi^*\approx0.50$ for $r_w=5$ Protocol~A.
  \item \textbf{FSS slope at $\xi=0.707$ is $-0.635$}, identifying
        $-\beta/\nu\approx-0.64$ and $\xi^*\approx0.65$--$0.71$.
        Implied $\nu\approx0.98$, consistent with Paper~63.
  \item \textbf{Indirect $\eta\approx1.28$} (via hyperscaling);
        exponent triple $(\beta,\nu,\eta)\approx(0.63,0.98,1.28)$ matches
        no known 2D non-equilibrium class.
  \item \textbf{Paper~72}: direct measurement of $\eta$ from the two-point
        fieldM correlator $G(r)\sim r^{-(d-2+\eta)}$ and check of the
        scaling relation $\beta=\nu\eta/2$ (in $d=2$).
\end{enumerate}

\end{document}
